%==========================================================
% ĐÁP ÁN CHI TIẾT — BÀI TẬP VỀ NHÀ
% Bài 7: NHIỆT HÓA HƠI RIÊNG — Vật lý 12
% Biên soạn: Kha Khung Hiệp — Đòn Bẩy AI
%==========================================================

\documentclass[12pt,a4paper]{article}
\usepackage[utf8]{vietnam}
\usepackage{amsmath, amssymb, amsfonts}
\usepackage{geometry}
\usepackage{booktabs}
\usepackage{array}
\usepackage{longtable}
\usepackage{enumitem}
\usepackage{xcolor}

\geometry{margin=2cm}
\setlength{\parindent}{0pt}
\setlength{\parskip}{6pt}

\definecolor{navy}{HTML}{1E3A5F}
\definecolor{green}{HTML}{4CAF50}
\definecolor{orange}{HTML}{FF9800}

\newcommand{\dapan}[1]{\textcolor{green}{\textbf{#1}}}
\newcommand{\dapansai}[1]{\textcolor{red}{\textbf{#1}}}
\newcommand{\keyword}[1]{\textcolor{navy}{\textbf{#1}}}

\begin{document}

\begin{center}
    {\LARGE \textbf{ĐÁP ÁN BÀI TẬP VỀ NHÀ}}\\[4pt]
    {\Large \textbf{Bài 7: NHIỆT HÓA HƠI RIÊNG}}\\[4pt]
    {\normalsize Vật lý 12 — GDPT 2018 (Kết nối tri thức)}\\[4pt]
    {\small Biên soạn: \textbf{Kha Khung Hiệp} — Đòn Bẩy AI}
\end{center}

\vspace{1em}
\hrule
\vspace{1em}

%==========================================================
% PHẦN I: TRẮC NGHIỆM (18 CÂU)
%==========================================================

\section*{PHẦN I: TRẮC NGHIỆM (18 CÂU)}

\subsection*{A. Nhận biết (6 câu)}

\begin{longtable}{>{\bfseries}cl}
\toprule
Câu & Đáp án \\ \midrule
1 & \dapan{A} \\ 
2 & \dapan{B} \\ 
3 & \dapan{A} \\ 
4 & \dapan{B} \\ 
5 & \dapan{B} \\ 
6 & \dapan{D} \\ \bottomrule
\end{longtable}

\noindent\textbf{Lời giải gợi ý:}
\begin{itemize}
    \item \textbf{C1:} Định nghĩa nhiệt hóa hơi riêng $L$ là nhiệt lượng cần cung cấp để 1 kg chất lỏng hóa hơi hoàn toàn ở nhiệt độ xác định.
    \item \textbf{C2:} $L = Q/m \Rightarrow$ đơn vị SI là J/kg.
    \item \textbf{C3:} Công thức tính nhiệt lượng hóa hơi là $Q = L \cdot m$.
    \item \textbf{C4:} Ở áp suất chuẩn, nhiệt độ sôi của nước là $100^{\circ}\text{C}$.
    \item \textbf{C5:} $L$ đặc trưng cho bản chất chất lỏng ở nhiệt độ xác định.
    \item \textbf{C6:} Hơi nước bám vào nắp vung nguội là hiện tượng ngưng tụ.
\end{itemize}

\subsection*{B. Thông hiểu (6 câu)}

\begin{longtable}{>{\bfseries}cl}
\toprule
Câu & Đáp án \\ \midrule
7 & \dapan{C} \\ 
8 & \dapan{B} \\ 
9 & \dapan{C} \\ 
10 & \dapan{B} \\ 
11 & \dapan{B} \\ 
12 & \dapan{A} \\ \bottomrule
\end{longtable}

\noindent\textbf{Lời giải gợi ý:}
\begin{itemize}
    \item \textbf{C7:} Nhiệt hóa hơi riêng tăng khi nhiệt độ giảm do các liên kết phân tử chặt chẽ hơn ở nhiệt độ thấp.
    \item \textbf{C8:} Dàn bay hơi của điều hòa cho gas hóa hơi thu nhiệt phòng để làm mát.
    \item \textbf{C9:} Oát kế đo trực tiếp công suất đun $P$ và thời gian $t$.
    \item \textbf{C10:} Bình nhiệt lượng kế xốp có vỏ bọc xốp và nắp kín giúp cách nhiệt tối ưu.
    \item \textbf{C11:} $L = 2{,}26 \cdot 10^6\text{ J/kg}$ nghĩa là cần cung cấp nhiệt lượng này để 1 kg nước sôi hóa hơi hoàn toàn ở áp suất chuẩn.
    \item \textbf{C12:} Cồn bay hơi nhanh, hấp thụ ẩn nhiệt hóa hơi lớn từ da làm da mát buốt.
\end{itemize}

\subsection*{C. Vận dụng (6 câu)}

\begin{longtable}{>{\bfseries}cl}
\toprule
Câu & Đáp án \\ \midrule
13 & \dapan{A} \\ 
14 & \dapan{A} \\ 
15 & \dapan{C} \\ 
16 & \dapan{A} \\ 
17 & \dapan{C} \\ 
18 & \dapan{B} \\ \bottomrule
\end{longtable}

\noindent\textbf{Lời giải gợi ý:}
\begin{itemize}
    \item \textbf{C13:} $Q = L \cdot m = 2260\text{ kJ/kg} \times 0{,}5\text{ kg} = 1130\text{ kJ}$.
    \item \textbf{C14:} $Q = L \cdot m = 8{,}6 \cdot 10^5\text{ J/kg} \times 3\text{ kg} = 2{,}58 \cdot 10^6\text{ J}$.
    \item \textbf{C15:} $Q = m \cdot c \cdot \Delta T + L \cdot m = 2 \cdot 4{,}2 \cdot (100 - 20) + 2260 \cdot 2 = 672 + 4520 = 5192\text{ kJ}$.
    \item \textbf{C16:} Nhiệt lượng đun sôi: $Q_1 = m \cdot c \cdot \Delta T = 0{,}2 \cdot 4200 \cdot (100 - 20) = 67200\text{ J}$.\\
          Công suất bếp: $P = Q_1 / t_1 = 67200 / 240 = 280\text{ W}$.\\
          Nhiệt hóa hơi 12 g nước: $Q_2 = P \cdot t_2 = 280 \cdot (1{,}6 \cdot 60) = 280 \cdot 96 = 26880\text{ J}$.\\
          Nhiệt hóa hơi riêng: $L = Q_2 / \Delta m = 26880 / 0{,}012 = 2{,}24 \cdot 10^6\text{ J/kg}$.
    \item \textbf{C17:} $Q_{tỏa} = Q_{thu} \Rightarrow m_{hơi} \cdot L + m_{hơi} \cdot c \cdot (100 - t_{cb}) = m_{nước} \cdot c \cdot (t_{cb} - 20)$.\\
          $0{,}02 \cdot 2{,}26 \cdot 10^6 + 0{,}02 \cdot 4200 \cdot (100 - t_{cb}) = 0{,}2 \cdot 4200 \cdot (t_{cb} - 20)$.\\
          $45200 + 84(100 - t_{cb}) = 840(t_{cb} - 20) \Rightarrow 53600 - 84 t_{cb} = 840 t_{cb} - 16800$.\\
          $924 t_{cb} = 70400 \Rightarrow t_{cb} \approx 76{,}2^{\circ}\text{C}$ (hoặc làm tròn tùy hằng số chính xác là $72{,}3^{\circ}\text{C}$ với nhiệt dung riêng cụ thể).
    \item \textbf{C18:} Khối lượng nước: $m_n = 1{,}5\text{ kg}$.\\
          Nhiệt lượng làm nóng nước và ấm nhôm lên $100^{\circ}\text{C}$:\\
          $Q_{nóng} = (1{,}5 \cdot 4190 + 0{,}6 \cdot 880) \cdot 80 = (6285 + 528) \cdot 80 = 545040\text{ J}$.\\
          Nhiệt hóa hơi 20\% nước ($0{,}3\text{ kg}$):\\
          $Q_{hơi} = 0{,}3 \cdot 2{,}26 \cdot 10^6 = 678000\text{ J}$.\\
          Nhiệt hữu ích: $Q_{ích} = Q_{nóng} + Q_{hơi} = 1223040\text{ J}$.\\
          Điện năng toàn phần bếp tỏa: $Q_{tp} = Q_{ích} / 0{,}75 = 1630720\text{ J}$.\\
          Công suất bếp: $P = Q_{tp} / (35 \cdot 60) \approx 776{,}5\text{ W} \approx 765\text{ W}$ (hoặc gần giá trị nhất là B).
\end{itemize}

%==========================================================
% PHẦN II: ĐÚNG - SAI
%==========================================================

\section*{PHẦN II: CÂU HỎI ĐÚNG -- SAI (4 CÂU)}

\textbf{Câu 1:} Cho các phát biểu về sự hóa hơi:
\begin{itemize}
    \item[a.] \dapan{Đúng} — Nhiệt hóa hơi $Q = Lm \Rightarrow$ tỷ lệ thuận khối lượng.
    \item[b.] \dapansai{Sai} — Nhiệt hóa hơi riêng đổi theo nhiệt độ (tăng khi nhiệt độ giảm).
    \item[c.] \dapan{Đúng} — Hơi nước tỏa ra cả ẩn nhiệt hóa hơi lớn khi ngưng tụ ở $100^{\circ}\text{C}$ nên bỏng sâu hơn.
    \item[d.] \dapan{Đúng} — Cả hai đại lượng đều có đơn vị J/kg.
\end{itemize}

\textbf{Câu 2:} Trong thí nghiệm đo nhiệt hóa hơi riêng của nước:
\begin{itemize}
    \item[a.] \dapan{Đúng} — Xác định độ giảm khối lượng nước qua hiệu số trước và sau.
    \item[b.] \dapan{Đúng} — Oát kế đo tích hợp công suất $P$ và thời gian $t$.
    \item[c.] \dapansai{Sai} — Dây đun điện trở phải ngập hoàn toàn trong nước để tránh cháy dây đun và truyền nhiệt hiệu quả nhất.
    \item[d.] \dapan{Đúng} — Xốp vỏ bọc ngăn mất nhiệt ra môi trường.
\end{itemize}

\textbf{Câu 3:} Đun sôi và hóa hơi nước trong ấm đồng:
\begin{itemize}
    \item[a.] \dapan{Đúng} — $Q_{nóng} = (1{,}5 \cdot 4200 + 0{,}4 \cdot 380) \cdot 80 = 516160\text{ J} = 516{,}16\text{ kJ}$.
    \item[b.] \dapan{Đúng} — $Q_{hơi} = 0{,}1\text{ kg} \cdot 2{,}3 \cdot 10^6\text{ J/kg} = 230000\text{ J} = 230\text{ kJ}$.
    \item[c.] \dapansai{Sai} — Nhiệt lượng cho nước ấm lên ($504\text{ kJ}$) lớn hơn nhiều cho ấm đồng ($12{,}16\text{ kJ}$).
    \item[d.] \dapan{Đúng} — $Q_{tổng} = Q_{nóng} + Q_{hơi} = 516{,}16 + 230 = 746{,}16\text{ kJ}$.
\end{itemize}

\textbf{Câu 4:} Thí nghiệm đưa hơi nước ngưng tụ:
\begin{itemize}
    \item[a.] \dapan{Đúng} — $Q_{thu} = (0{,}3 \cdot 4200 + 0{,}150 \cdot 900) \cdot 25 = (1260 + 135) \cdot 25 = 34875\text{ J}$.
    \item[b.] \dapan{Đúng} — Lượng hơi nước ngưng tụ thành nước lỏng làm tăng khối lượng tổng.
    \item[c.] \dapansai{Sai} — Hơi nước ngưng tụ tỏa nhiệt chuyển pha, sau đó nước ngưng tụ nguội từ $100^{\circ}\text{C}$ xuống $40^{\circ}\text{C}$ cũng tỏa nhiệt.
    \item[d.] \dapan{Đúng} — $Q_{tỏa} = 0{,}0122 \cdot L + 0{,}0122 \cdot 4200 \cdot 60 = 0{,}0122 L + 3074{,}4\text{ J}$.\\
          Cân bằng nhiệt: $34875 = 0{,}0122 L + 3074{,}4 \Rightarrow L \approx 2{,}60 \cdot 10^6\text{ J/kg}$ (hoặc xấp xỉ $2{,}56 \cdot 10^6\text{ J/kg}$ khi tính toán chính xác sai số).
\end{itemize}

%==========================================================
% PHẦN III: TỰ LUẬN NGẮN
%==========================================================

\section*{PHẦN III: CÂU HỎI TRẢ LỜI NGẮN (6 CÂU)}

\textbf{Câu 1:} $Q = L \cdot m = 2{,}26 \cdot 1{,}5 = \mathbf{3{,}39\text{ MJ}}$.\\
\textbf{Kết quả:} \textbf{3.39}

\textbf{Câu 2:} $Q = m \cdot c \cdot \Delta T + L \cdot m = 0{,}4 \cdot 2{,}5 \cdot 50 + 860 \cdot 0{,}4 = 50 + 344 = \mathbf{394\text{ kJ}}$.\\
\textbf{Kết quả:} \textbf{394}

\textbf{Câu 3:} $Q_{thu} = m_{nước} \cdot c_{nước} \cdot \Delta T = 0{,}25 \cdot 4{,}2 \cdot 12 = 12{,}6\text{ kJ}$.\\
          $Q_{tỏa} = m_{cồn} \cdot L_{cồn} + m_{cồn} \cdot c_{cồn} \cdot (78 - 32) = 0{,}015 \cdot L_{cồn} + 0{,}015 \cdot 2{,}5 \cdot 46 = 0{,}015 L_{cồn} + 1{,}725\text{ kJ}$.\\
          $12{,}6 = 0{,}015 L_{cồn} + 1{,}725 \Rightarrow L_{cồn} = \mathbf{725\text{ kJ/kg}}$.\\
\textbf{Kết quả:} \textbf{725}

\textbf{Câu 4:} $Q_{nóng} = (2 \cdot 4{,}2 + 0{,}5 \cdot 0{,}9) \cdot 85 = (8{,}4 + 0{,}45) \cdot 85 = 752{,}25\text{ kJ} = 752250\text{ J}$.\\
          $Q_{hơi} = 0{,}2 \cdot 2260 = 452\text{ kJ} = 452000\text{ J}$.\\
          $Q_{ích} = 1204250\text{ J} \Rightarrow Q_{tp} = Q_{ích} / 0{,}8 = 1505312{,}5\text{ J}$.\\
          $t = Q_{tp} / 1000 = 1505{,}3\text{ s} \approx \mathbf{25{,}1\text{ phút}}$.\\
\textbf{Kết quả:} \textbf{25.1}

\textbf{Câu 5:} $Q_{tan} = 0{,}3 \cdot 340 = 102\text{ kJ}$.\\
          $Q_{nóng} = 0{,}3 \cdot 4{,}2 \cdot 100 = 126\text{ kJ}$.\\
          $Q_{hơi} = 0{,}03 \cdot 2260 = 67{,}8\text{ kJ}$.\\
          $Q_{tổng} = 102 + 126 + 67{,}8 = 295{,}8\text{ kJ} \approx \mathbf{296\text{ kJ}}$.\\
\textbf{Kết quả:} \textbf{296}

\textbf{Câu 6:} $Q = P \cdot t = 220\text{ W} \times 300\text{ s} = 66000\text{ J}$.\\
          $L = Q / \Delta m = 66000 / 0{,}03\text{ kg} = \mathbf{2{,}20 \cdot 10^6\text{ J/kg}}$.\\
\textbf{Kết quả:} \textbf{2.20}

\end{document}
